\documentclass[12pt]{article}

\usepackage{amsmath,amssymb,amsthm}
\usepackage{geometry}
\usepackage[colorlinks]{hyperref}

\geometry{margin=1in}

\title{Formal Appendix: State--Collapse Theorem}
\author{Jinu Jang}
\date{v2.0 --- corrected; every statement below is machine-checked in Lean~4}

\theoremstyle{plain}
\newtheorem{theorem}{Theorem}
\newtheorem{lemma}{Lemma}
\theoremstyle{remark}
\newtheorem*{remark}{Remark}

% shorthand
\newcommand{\Set}{\mathbf{Set}}
\newcommand{\Vect}{\mathbf{Vect}_{\mathbb R}}
\newcommand{\id}{\operatorname{id}}
\newcommand{\RR}{\mathbb R}
\newcommand{\Kl}{\operatorname{Kl}}

\begin{document}
\maketitle

\begin{abstract}
This appendix states and proves the sense in which the multiplication of the
state monad is an idempotent projection. Sections~0--4 and~7 are essentially
unchanged from v1.1. Sections~5 and~6 are \emph{corrected}: the reshuffle
isomorphism asserted there does not exist, and \S\ref{sec:noreshuffle} gives a
counterexample. Section~9 is an erratum listing the changes, and \S10 records
what the theorem does \emph{not} show. Lean names accompanying each statement
refer to the formalisation at
\url{https://github.com/Kairose-master/mu_eq_pi}.
\end{abstract}

\section*{0.\;Notation}
\begin{itemize}
  \item $\Set$ -- category of sets and functions.
  \item $\Vect$ -- category of real vector spaces and linear maps.
  \item $S$ -- fixed non-empty \emph{state set}.
  \item $\RR^{X}$ -- free vector space on $X$ with basis
        $\{\delta_x\}_{x\in X}$.
  \item $\otimes$ -- algebraic tensor product over $\RR$.
  \item $\RR^{f}$ -- the linear map $\delta_a\mapsto\delta_{f(a)}$ induced by a
        function $f$.
\end{itemize}

\bigskip
\section{The state monad \texorpdfstring{$T_S$}{T\_S}}
\[
T_S X \;:=\; S \Rightarrow (S\times X).
\]

\paragraph{Unit.}
\[
\eta_X(x)(s)\;:=\;(s,x).
\]

\paragraph{Multiplication.}
For $\Phi\in T_ST_SX$, write $\Phi(s)=(s',\psi)$ with $s'\in S$ and
$\psi\in T_SX$. Then
\[
\mu_X : T_S T_S X \longrightarrow T_S X,
\qquad
\mu_X(\Phi)(s) \;:=\; \bigl(\pi_2\Phi(s)\bigr)\bigl(\pi_1\Phi(s)\bigr)
\;=\;\psi(s').
\]

\paragraph{Monad laws.}
\[
\mu_X \circ T_S\mu_X = \mu_X \circ \mu_{T_S X},
\qquad
\mu_X\circ\eta_{T_SX}=\id,
\qquad
\mu_X \circ T_S\eta_X \;=\; \id .
\tag{$*$}
\]
\emph{Lean:} \verb|St.mu_assoc|, \verb|St.mu_eta|, \verb|St.mu_map_eta|.

\begin{remark}
The right unit law in $(*)$ is the engine of everything below: it says $\mu_X$
is a \emph{split epimorphism} with canonical section $T_S\eta_X$. Note that
``$\mu\circ\mu=\mu$'' is not a statement one can make --- $\mu$ has different
source and target --- whereas $T_S\eta\circ\mu$ is an endomorphism of $T_S^2$
and is idempotent. \emph{Lean:} \verb|St.collapse_idem|.
\end{remark}

\bigskip
\section{Free \texorpdfstring{$T_S$}{T\_S}-algebras}
A $T_S$-algebra is a pair $(A,\alpha)$ with $\alpha:T_SA\to A$ obeying the usual
axioms. The free algebra on a set $X$ is $\bigl(T_S X,\;\mu_X\bigr)$, and
$\operatorname{Alg}_{T_S}$ has morphisms $f:(A,\alpha)\to(B,\beta)$ satisfying
$\beta\circ T_S f = f\circ\alpha$.

\bigskip
\section{The vector-space functor \texorpdfstring{$\mathcal F$}{F}}
\label{sec:functor}

On objects, put
\[
\mathcal F(X)\;=\;\RR^{S}\otimes\RR^{X}\;\cong\;\RR^{\,S\times X}.
\]
The set $C_X:=S\times X$ is the set of \emph{configurations}: a state together
with a value.

On morphisms, $\mathcal F$ must be defined on \emph{Kleisli} arrows, not on
functions $X\to Y$. Uncurrying gives a bijection
\[
\Kl(T_S)(X,Y)\;=\;\bigl(X\to S\Rightarrow(S\times Y)\bigr)
\;\cong\;\bigl(C_X\to C_Y\bigr),
\qquad
k\;\longmapsto\;\widehat k,\quad \widehat k(s,x)=k(x)(s),
\]
and this bijection carries Kleisli composition to ordinary composition of
functions. Setting
\[
\boxed{\;\mathcal F(k)\;:=\;\RR^{\widehat k}\;:\;\RR^{C_X}\to\RR^{C_Y}\;}
\]
therefore defines a functor $\mathcal F:\Kl(T_S)\to\Vect$.
\emph{Lean:} \verb|threadEquiv|, \verb|uncur|, \verb|Lin|.

\begin{lemma}
$\mathcal F$ is faithful but not full. Distinct Kleisli arrows have distinct
matrices; conversely the image of $\mathcal F$ consists of the $0$--$1$ matrices
with exactly one $1$ in each column, so e.g.\ $2\cdot\id$ is not of the form
$\mathcal F(k)$.
\end{lemma}
\emph{Lean:} \verb|instFaithfulLin|, \verb|two_smul_id_not_lin|.

\begin{remark}
$\Kl(T_S)$ is, by the displayed bijection, the co-Kleisli category of the
comonad $-\times S$. Identifying the two is what makes a linear model possible:
a Kleisli arrow $X\to S\Rightarrow(S\times Y)$ is higher-order and does not
embed in a vector space, but $\widehat k$ is a function between plain sets, and
functions between sets linearise.
\end{remark}

\bigskip
\section{Split idempotent pair \texorpdfstring{$P,\,i_\sigma$}{P, i}}
\label{sec:split}

Let $D_X:=S\times C_X$ be the \emph{doubled configuration set}, so that
\[
\RR^{D_X}\;\cong\;(\RR^{S}\otimes\RR^{S})\otimes\RR^{X},
\]
with basis $e_{s_1}\otimes e_{s_2}\otimes\delta_x$. Here $s_1$ is the
\emph{inner} slot --- the state at which a second stage is to be started, a free
parameter before the two stages are joined --- and $s_2$ is the \emph{outer}
slot, part of the incoming configuration.

\paragraph{The projection.}
\[
P_{S,X}:(\RR^{S}\otimes\RR^{S})\otimes\RR^{X}\to \RR^{S}\otimes\RR^{X},
\qquad
P\bigl(e_{s_1}\otimes e_{s_2}\otimes\delta_x\bigr)=e_{s_2}\otimes\delta_x .
\]

\paragraph{The section.}
For a \emph{transition function} $\sigma:C_X\to S$ put
\[
i_\sigma:\RR^{S}\otimes\RR^{X}\to(\RR^{S}\otimes\RR^{S})\otimes\RR^{X},
\qquad
i_\sigma\bigl(e_s\otimes\delta_x\bigr)=e_{\sigma(s,x)}\otimes e_{s}\otimes\delta_x .
\]
Version~1.1 wrote only the special case $\sigma(s,x)=s$, giving the diagonal
$i(e_s\otimes\delta_x)=(e_s\otimes e_s)\otimes\delta_x$; the general theorem
needs the index $\sigma$.

\begin{lemma}[Retraction]
$P\circ i_\sigma=\id_{\RR^{S}\otimes\RR^{X}}$ for every $\sigma$, and hence
\[
e_\sigma:=i_\sigma\circ P
\qquad\text{satisfies}\qquad
e_\sigma^{\,2}=e_\sigma .
\]
On the basis, $e_\sigma\bigl(e_{s_1}\otimes e_{s_2}\otimes\delta_x\bigr)
= e_{\sigma(s_2,x)}\otimes e_{s_2}\otimes\delta_x$: the free inner slot is
overwritten by the value $\sigma$ prescribes. Thus $e_\sigma$ is the projection
of the doubled state space onto the \emph{graph of $\sigma$}.
\end{lemma}
\emph{Lean:} \verb|Flat_comp_Sect|, \verb|proj_comp_proj|, \verb|proj_single|.

\begin{lemma}[Properness]
If $|S|\ge 2$ and $X\neq\emptyset$ then $e_\sigma\neq\id$. Moreover, for finite
$S$ and $X$,
\[
\operatorname{rank}e_\sigma=\operatorname{tr}e_\sigma=|S|\,|X|,
\qquad
\dim\ker e_\sigma=|S|\,|X|\,(|S|-1),
\]
independently of $\sigma$.
\end{lemma}
\emph{Lean:} \verb|proj_ne_id|, \verb|finrank_range_proj|, \verb|trace_proj|,
\verb|finrank_ker_proj|.

\paragraph{Why $P$ is well defined.}
$P$ is \emph{not} an instance of a projection $V\otimes W\to W$: no such natural
map exists for general modules. What makes $P$ legitimate is that $\RR^{S}$ is
\emph{free} on $S$, hence carries an augmentation
\[
\varepsilon:\RR^{S}\to\RR,\qquad \varepsilon(e_s)=1,
\]
and $P=\lambda\circ(\varepsilon\otimes\id\otimes\id)$ where $\lambda$ is the
unitor $\RR\otimes V\cong V$. \emph{Lean:} \verb|aug|, \verb|Flat_eq_aug|.

\paragraph{Naturality.}
For every function $f:X\to Y$, $P_{S,Y}\circ(\id\otimes\id\otimes\RR^{f})
=(\id\otimes\RR^{f})\circ P_{S,X}$, and likewise for $i_\sigma$ and $e_\sigma$
once $\sigma$ is transported along $f$.

\bigskip
\section{\texorpdfstring{$\mathcal F(T_SX)$}{F(T\_S X)} admits no reshuffle}
\label{sec:noreshuffle}

Version~1.1 asserted an isomorphism
\[
\theta_X:\mathcal F(T_S X)\;\xrightarrow{\ \cong\ }\;
(\RR^{S}\otimes\RR^{S})\otimes\RR^{X},
\qquad
\theta_X\bigl(e_{s_1}\otimes\delta_{\langle s_2,x\rangle}\bigr)
=e_{s_1}\otimes e_{s_2}\otimes\delta_x .
\]
\textbf{This map does not exist.} The formula indexes basis vectors of
$\RR^{T_SX}$ by pairs $\langle s_2,x\rangle\in S\times X$, i.e.\ it silently
replaces $T_SX=S\Rightarrow(S\times X)$ by $S\times X$. These are different
sets, and no isomorphism of the stated shape exists:

\begin{theorem}[Counterexample]
Let $|S|=2$ and $|X|=1$. Then $|T_SX|=|S\times X|^{|S|}=2^{2}=4$, so
\[
\dim\mathcal F(T_SX)=|S|\cdot|T_SX|=2\cdot 4=8,
\qquad
\dim\bigl((\RR^{S}\otimes\RR^{S})\otimes\RR^{X}\bigr)=2\cdot 2\cdot 1=4 .
\]
Since $8\neq4$, no linear isomorphism $\theta_X$ exists.
\end{theorem}

In general $\dim\mathcal F(T_SX)=|S|\cdot(|S||X|)^{|S|}$, which grows
exponentially in $|S|$, against $|S|^2|X|$ on the right.

\paragraph{The correct route.}
Do not apply $\mathcal F$ to the \emph{object} $T_SX$. The object assignment of
\S\ref{sec:functor} is already correct; what must be taken on Kleisli arrows is
the \emph{morphism} assignment. The doubled configuration set $D_X=S\times C_X$
of \S\ref{sec:split} is a plain finite set, and it --- not $\mathcal F(T_SX)$
--- is where $P$ and $i_\sigma$ live.

\bigskip
\section{\texorpdfstring{$\mathcal F(\mu)$}{F(mu)} explicitly}

Fix a nested Kleisli arrow $t:X\to T_ST_SY$. Writing
$t(x)(s)=(s',\psi)$, define
\[
\sigma_t:C_X\to S,\quad \sigma_t(s,x):=s'=\pi_1\bigl(t(x)(s)\bigr),
\]
\[
\widetilde t:D_X\to C_Y,\quad
\widetilde t\bigl(s_1,(s,x)\bigr):=\psi(s_1)=\bigl(\pi_2 (t(x)(s))\bigr)(s_1),
\]
and $\gamma_\sigma:C_X\to D_X$, $\gamma_\sigma(p):=(\sigma(p),p)$, so that
$i_\sigma=\RR^{\gamma_\sigma}$ and $P=\RR^{\pi_2}$.

\begin{theorem}[Flattening is substitution]
\label{thm:subst}
\[
\boxed{\;\widehat{\mu\circ t}\;=\;\widetilde t\circ\gamma_{\sigma_t}\;}
\qquad\text{hence}\qquad
\mathcal F(\mu\circ t)\;=\;\RR^{\widetilde t}\circ i_{\sigma_t}.
\]
\end{theorem}

\begin{proof}
Both sides at $(s,x)$: the left is $\mu(t(x))(s)=\psi(s')$ by \S1, and the
right is $\widetilde t(\sigma_t(s,x),(s,x))=\widetilde t(s',(s,x))=\psi(s')$.
Apply $\RR^{(-)}$, which is functorial.
\end{proof}

Thus flattening is exactly \emph{substitution of the transition function into
the free inner slot}. \emph{Lean:} \verb|threadEquiv_mu|, \verb|lin_mu|.

\bigskip
\section{Collapse Theorem}

\begin{theorem}[State--Collapse]
For every $\sigma:C_X\to S$ the operator $e_\sigma=i_\sigma\circ P$ is a split
idempotent in $\Vect$, and for every nested Kleisli arrow $t:X\to T_ST_SY$ the
square
\[
\boxed{\;\RR^{\widetilde t}\circ e_{\sigma_t}
\;=\;\mathcal F(\mu\circ t)\circ P\;}
\]
commutes: projecting the doubled state space onto the graph of $\sigma_t$ and
then running is the same as flattening and then running.
\end{theorem}

\begin{proof}
Idempotence: $e_\sigma^{2}=i_\sigma\circ(P\circ i_\sigma)\circ P
=i_\sigma\circ P=e_\sigma$ by the Retraction Lemma. For the square,
\[
\RR^{\widetilde t}\circ e_{\sigma_t}
=\RR^{\widetilde t}\circ i_{\sigma_t}\circ P
=\mathcal F(\mu\circ t)\circ P
\]
by Theorem~\ref{thm:subst}.
\end{proof}
\emph{Lean:} \verb|mu_eq_pi|.

\bigskip
\section{Lean identifiers}

The formalisation orders configurations as $X\times S$ rather than $S\times X$
and places the free slot second, so $P$ appears as $\RR^{\pi_1}$ there.

\begin{center}
\renewcommand{\arraystretch}{1.25}
\begin{tabular}{@{}lll@{}}
\textbf{Symbol} & \textbf{Lean name} & \textbf{File}\\\hline
$\mu_X$ & \verb|St.mu| & \verb|StateMonad.lean|\\
$\mu\circ T_S\eta=\id$ & \verb|St.mu_map_eta| & \verb|StateMonad.lean|\\
$T_S\eta\circ\mu$ idempotent & \verb|St.collapse_idem| & \verb|StateMonad.lean|\\
$\widehat{(-)}$ & \verb|threadEquiv| & \verb|Kleisli.lean|\\
$\mathcal F$ & \verb|Lin| & \verb|Linearization.lean|\\
faithful, not full & \verb|instFaithfulLin|, \verb|two_smul_id_not_lin| & \verb|Linearization.lean|\\
$P_{S,X}$ & \verb|Flat| & \verb|Projection.lean|\\
$i_\sigma$ & \verb|Sect| & \verb|Projection.lean|\\
$e_\sigma$ & \verb|proj| & \verb|Projection.lean|\\
$P\circ i_\sigma=\id$ & \verb|Flat_comp_Sect| & \verb|Projection.lean|\\
$e_\sigma^2=e_\sigma$ & \verb|proj_comp_proj| & \verb|Projection.lean|\\
$e_\sigma\neq\id$ & \verb|proj_ne_id| & \verb|Projection.lean|\\
$P=\lambda\circ(\varepsilon\otimes\id)$ & \verb|Flat_eq_aug| & \verb|Projection.lean|\\
rank, trace, kernel & \verb|finrank_range_proj|, \verb|trace_proj| & \verb|Projection.lean|\\
$\sigma_t$, $\widetilde t$ & \verb|outerNext|, \verb|expand| & \verb|Collapse.lean|\\
Theorem~\ref{thm:subst} & \verb|threadEquiv_mu|, \verb|lin_mu| & \verb|Collapse.lean|\\
Collapse Theorem & \verb|mu_eq_pi| & \verb|Collapse.lean|\\
\end{tabular}
\end{center}

Every entry compiles against Lean~4 (\verb|v4.32.1|) and \verb|mathlib|, with no
\verb|sorry| and no axioms beyond \verb|propext|, \verb|Classical.choice| and
\verb|Quot.sound|.

\bigskip
\section{Erratum: changes from v1.1}

\begin{enumerate}
  \item \S1, multiplication: the formula read
        $\mu_X(\Phi)(s)=\Phi(s)(\pi_1\Phi(s))$, applying a \emph{pair} to an
        argument. Corrected to
        $\mu_X(\Phi)(s)=(\pi_2\Phi(s))(\pi_1\Phi(s))$.
  \item \S3: $\mathcal F$ was given only on functions $X\to Y$. To be a functor
        on $\Kl(T_S)$ it must be defined on Kleisli arrows, via uncurrying.
  \item \S4: the section was stated only in the state-preserving case. The
        general statement requires the family $i_\sigma$ indexed by a transition
        function. The justification of $P$ via the augmentation of a free module
        has been added, since $P$ is \emph{not} an instance of a nonexistent
        projection $V\otimes W\to W$.
  \item \S5: the reshuffle isomorphism $\theta_X$ does not exist; see the
        counterexample in \S\ref{sec:noreshuffle}. This was the substantive
        error of v1.1.
  \item \S6: consequently $\mathcal F(\mu_X)=\theta^{-1}\circ P\circ\theta$ was
        ill-typed. Replaced by Theorem~\ref{thm:subst}, which factors
        $\mu\circ t$ for each nested Kleisli arrow $t$ rather than trying to
        linearise the object $T_SX$.
  \item \S8: the Lean identifiers listed
        (\verb|state_drop|, \verb|state_dup|, \verb|state_idem|,
        \verb|collapse_split|, \verb|collapse_idem|, \verb|collapse_theorem|)
        did not exist in the repository. Replaced by the actual names.
\end{enumerate}

\bigskip
\section{What the theorem does not show}

\begin{enumerate}
  \item \textbf{Not specific to the state monad.} That $T\eta\circ\mu$ is
        idempotent is a restatement of the right unit law and holds for
        \emph{every} lawful monad. \emph{Lean:} \verb|mcollapse_idem|,
        \verb|St.collapse_eq_mcollapse|.
  \item \textbf{Not specific to linear algebra.} $e_\sigma$ is exactly the
        linearisation of the idempotent \emph{set map}
        $(s_1,p)\mapsto(\sigma(p),p)$, whose idempotence is immediate. Its
        matrix is a $0$--$1$ matrix with one $1$ per column, its spectrum lies
        in $\{0,1\}$ for the trivial reason, and $\operatorname{rank}=
        \operatorname{tr}=|S||X|$ is a count of fixed points of that set map.
        The vector space is bookkeeping. \emph{Lean:}
        \verb|proj_eq_lin_setIdem|, \verb|fixedPoints_setIdem|.
  \item \textbf{The projection depends on the computation.} There is no single
        $e$ doing the work of $\mu$ for all nested arrows at once; only the rank
        and trace are independent of $\sigma$.
  \item \textbf{No claim is made} about language models, ``tensor collapse'', or
        any domain-specific language.
\end{enumerate}

A monad whose linearisation is \emph{not} a $0$--$1$ matrix --- the distribution
monad, whose Kleisli arrows become stochastic matrices --- is where a linear
model would begin to carry content of its own.

\end{document}
